\begin{papercard}{wharton1992}
\cf{Stages} \stagetag{W}. An experience report on the \emph{second}, form-based cognitive walkthrough applied by developer teams to three complex systems. It is the source of most of the practical rules on task selection, granularity, group process and interpreting results that later versions adopted.

\cf{Aim and evidence type} To report how the CW worked when small teams of professional developers critiqued three complex systems, organised around five issues: task selection and coverage, the process, evaluator knowledge, group walkthroughs, and interpreting results \pc{wharton1992}{1}{We focus on five core issues: (1) task selection, coverage, and evaluation}. The authors state explicitly that the evidence is anecdotal and not experimental, and argue that recurrence across independent applications lends robustness \pc{wharton1992}{1}{one must keep in mind that we are reporting anecdotal evidence, not the results of controlled experiments}.

\cf{The method as used} A task-based method. Each step is judged by whether the typical user's goals and knowledge would lead to the action that brings them closer to the goal, and both problematic and successful steps are recorded \pc{wharton1992}{2}{Steps are deemed successful if the expected goals and knowledge of the typical user would result in selection}. Three phases: preparation (tasks, user information), evaluation (questions per step), and an \emph{ad hoc} interpretation phase \pc{wharton1992}{2}{The Walkthrough method consists of three basic phases}. The per-action form (Figure 1) asks about action availability, identifier location and wording, the link between identifier and action and between identifier and goal, each with an estimate of how many users would miss it (0--100\,\%) \pc{wharton1992}{2}{Is it obvious to tbe user that this action is a possible choice here}.

\cf{Cases} HP-VUE (a visual Unix desktop, beta version): three software engineers, seven tasks chosen by a CW expert rather than by the evaluators \pc{wharton1992}{2}{The evaluation of this interface was done by a group of three software engineers}. REPS (a sales-call system taken down after complaints): a developer, a requirements analyst and a cognitive psychologist; three tasks with forms and a fourth answered informally \pc{wharton1992}{3}{In a fourth task scenario, no forms were used}. BCA (a CAD prototype for circuit boards): a six-person design team, two tasks, one simple and one complex \pc{wharton1992}{3}{Two tasks were evaluated, one simple and the other complex}.

\cf{Task selection and coverage} CW gives no guidance on choosing tasks, because task selection lies outside its theory \pc{wharton1992}{3}{the Walkthrough methodology does not provide guidance on how to select tasks}. Tasks that mirror the designers' functional decomposition expose few problems. Choose realistic tasks that combine several core functions, so transitions between subtasks are covered, and balance realism against the number of tasks \pc{wharton1992}{3}{it is most important to choose realistic tasks which exercise key system functionality}. Boundaries need a decision: should tasks cover what the application does now or what users expect it to do? \pc{wharton1992}{4}{Should the tasks to be evaluated encompass only what the application currently does or what users will expect it to do}. Coverage was small: evaluators stopped after 6--15 hours and 2--7 tasks, while a full evaluation would take over 100 hours \pc{wharton1992}{4}{Because all evaluators stopped after 6-15 total hours and 2-7 tasks}. Only a single, ``most obvious'' path was evaluated per task \pc{wharton1992}{4}{We have always evaluated only a single path assuming that this is the one the user would choose}.

\cf{Granularity and repeated subtasks} Action granularity must be fixed for the interface: a set of keystrokes such as a file name can count as one action, but different physical actions (keystroke vs.\ mouse click) are separate \pc{wharton1992}{4}{In most cases it seems that a reasonable collection of keystrokes}. Identical subtasks in different tasks can behave differently, so all transitions must be evaluated, though some repeated subtask evaluations may be short-circuited \pc{wharton1992}{4}{for identical embedded actions, all transitions definitely should be evaluated}. CW has no high-level view of whether a task as a whole is sensible, so informal critiques were added \pc{wharton1992}{4}{We have relied on informal critiques to ensure that such issues are not missed}.

\cf{Process problems and fixes} Repeated form filling was so tedious that some evaluators were discouraged from using the method again \pc{wharton1992}{4}{found the repetitive form filling to be very tedous}. Side issues not relevant to the current task need an explicit tracking procedure with enough context to reconstruct them later \pc{wharton1992}{5}{the methodology needs explicit procedures for keeping track of side issues}. Developer-evaluators wanted to design and re-walk fixes immediately. The looser ``broad brush'' application on REPS's final task was more successful and satisfying, but only after structured practice \pc{wharton1992}{5}{We found the looser application of the method, as done by the REPS group on their final task}. Proposed aids: group-visible flip charts for user assumptions, rotating record-keeping, and a ``review card'' summarising the questions, with only problems committed to paper \pc{wharton1992}{5}{they would have preferred to have a}.

\cf{Evaluator knowledge} The method presupposes cognitive-science concepts (goal structures, supergoal kill-off) that developers lack \pc{wharton1992}{5}{The Walkthrough methodology presupposes more knowledge of cognitive science terms}. The commonest error was writing goals as interface operations (``click'', ``wait''). This hid a real problem, the missing feedback during a long file load \pc{wharton1992}{6}{This problem was not noted by the evaluators because they had assumed that the user had the goal of waiting}. Recommendation: include someone with at least moderate HCI or cognitive-science background \pc{wharton1992}{6}{A better approach would be to enlist someone with at least a moderate HCI}.

\cf{Group walkthroughs} Teams of two to six. Success correlated with division of duties, team size, session length, and the length and format of forms \pc{wharton1992}{6}{These are: division of duties, size of the team, length of the sessions}. A ``walkthrough advocate'' acted as facilitator, balancing discussion against refocusing \pc{wharton1992}{6}{The key issue for leading a group walkthrough seems to be the flexible managing}. Six people were too many and three were suggested \pc{wharton1992}{6}{The BCA evaluators felt that having six people on the walkthrough team was too many}. Long sessions (six hours) hurt more than breaks between two-hour sessions \pc{wharton1992}{6}{tiring the team with lengthy meetings seemed more damaging to the process}. Side benefits: learning about users and task analysis, and a neutral ground between requirements and development staff.

\cf{Interpreting results} The CW identifies mismatches between system affordances and user goals rather than problems as such. It catches linguistic mismatches (labels, menu items) better than graphical ones \pc{wharton1992}{7}{Walkthrough seems to do a better job of pointing out mismatches that are linguistic in origin}. Step focus produced suboptimal fixes (highlighting the next field instead of reorganising the screen) and erroneous ones that removed functionality needed by other tasks \pc{wharton1992}{7}{one evaluation team proposed a change that would produce a simpler action sequence}. These ``set effects'' call for global methods as a supplement \pc{wharton1992}{7}{methods that assess an interface more globally are needed as a supplement}.

\cf{Conclusions and recommendations (Figure 2)} Not ready for real teams ``without substantial extensions'' \pc{wharton1992}{7}{Based on our experiences, we say}. Figure 2: start simple, then include at least one task of typical complexity; choose realistic multi-function tasks; decide granularity, identical subtasks and variants; have one member who knows the theory; provide tracking aids; minimise bookkeeping; manage the group; and break sessions of 2--3 hours only between tasks \pc{wharton1992}{7}{Break the walkthrough up into sessions of reasonable leugth}. Remaining limits: poor fit with development practice, task selection (``the task suite should be generated by a knowledgeable professional''), and the low-level focus \pc{wharton1992}{8}{the task suite should be generated by a knowledgeable professional}.

\cf{Use for this thesis} Practical rules for the walkthrough protocol. (1) Justify the task set in the protocol, because CW gives none. Use realistic privacy tasks that cross screens (store page to onboarding to settings), not one task per menu item. (2) Fix and record the action granularity and the treatment of repeated subtasks, such as the same consent dialogue reached by two tasks. (3) Evaluate one documented correct path and log alternative paths as side issues. (4) Keep a side-issue log with screen context. (5) Report findings as goal--affordance \emph{mismatches} and keep proposed fixes separate. (6) Complement step-level judgements with a global pass, for example heuristics or a whole-task summary. The 6--15-hour coverage figure supports a small, justified task set per app.
\end{papercard}
